% ============================================================================
% BT#11 — Supplemental Material
%
% "Supplemental Material for: A Conditional Landauer–Nelson Bridge
% Supplemental Material for BT#11
% SHA-256: 781161D1440DA97C293630247D2B719C9673C0C2761799B0F763C93AD67707FB
% Lock date: 2026-03-18
% ============================================================================

\documentclass[twocolumn,showpacs,preprintnumbers,amsmath,amssymb]{revtex4-2}

\usepackage{amsmath,amssymb,amsfonts}
\usepackage{graphicx}
\usepackage{hyperref}
\usepackage{bm}
\usepackage{longtable}

\newcommand{\kB}{k_{\mathrm{B}}}
\newcommand{\Dth}{D_{\mathrm{th}}}
\newcommand{\Dq}{D_{\mathrm{q}}}
\newcommand{\geff}{\gamma_{\mathrm{eff}}}
\newcommand{\tcor}{\tau_{\mathrm{cor}}}
\newcommand{\Ebr}{E_{\mathrm{bridge}}}
\newcommand{\EL}{E_{\mathrm{Landauer}}}
\newcommand{\CTICE}{C_{\mathrm{TICE}}}

\begin{document}

\title{Supplemental Material for:\\
A Conditional Landauer--Nelson Bridge\\
with Non-Markovian TICE Closure}

\author{Kevin Henry Miller}
\email{Kevin@qbondnetwork.com}
\homepage{https://www.qbondnetwork.com}
\affiliation{Founder and President, Q-Bond Network DeSCI DAO, LLC}
\noaffiliation
% ORCID: 0009-0007-7286-3373

\date{\today}

\maketitle

% ============================================================================
\section{S1. Full Derivation Chain}
\label{sec:derivation}
% ============================================================================

\subsection{S1.1. Caldeira--Leggett: bath spectral density to noise kernel}

Start from the Ohmic--Drude spectral density:
\begin{equation}
  J(\omega) = m\gamma_0 \omega \,
  \frac{\Omega^2}{\omega^2 + \Omega^2} .
  \label{eq:spectral}
\end{equation}
The symmetrized noise kernel at temperature $T$ is:
\begin{equation}
  N(\tau) = \frac{1}{\pi}\int_0^\infty
  J(\omega)\coth\!\Bigl(\frac{\hbar\omega}{2\kB T}\Bigr)
  \cos(\omega\tau)\,d\omega .
  \label{eq:Nfull}
\end{equation}
In the high-$T$ limit $\hbar\omega \ll \kB T$ for all
$\omega$ in the support of~\eqref{eq:spectral}:
\begin{equation}
  \coth\!\Bigl(\frac{\hbar\omega}{2\kB T}\Bigr)
  \approx \frac{2\kB T}{\hbar\omega} ,
\end{equation}
so that
\begin{equation}
  N(\tau) \simeq \frac{2m\gamma_0\kB T\Omega^2}{\pi\hbar}
  \int_0^\infty
  \frac{\cos(\omega\tau)}{\omega^2 + \Omega^2}\,d\omega
  = \frac{m\gamma_0\kB T\Omega}{\hbar}\,e^{-\Omega|\tau|} .
  \label{eq:Nhigh}
\end{equation}

\subsection{S1.2. KMS condition to fluctuation--dissipation relation}

The KMS (Kubo--Martin--Schwinger) condition for the bath
correlation function at inverse temperature $\beta = 1/(\kB T)$
gives:
\begin{equation}
  \tilde{N}(\omega)
  = \coth\!\Bigl(\frac{\hbar\omega}{2\kB T}\Bigr)\,
  \tilde{D}(\omega) ,
  \label{eq:KMS}
\end{equation}
where $\tilde{D}(\omega) = J(\omega)/\pi$.
In the high-$T$ Markov limit this reduces to:
\begin{equation}
  N(\tau) = \frac{\kB T}{\hbar}\,\eta(\tau) ,
  \label{eq:FDT_highT}
\end{equation}
where $\eta(\tau) = m\gamma_0\,\Omega\,e^{-\Omega|\tau|}$ is
the dissipation kernel.

\subsection{S1.3. Green--Kubo: friction to thermal diffusion}

From the generalized Langevin equation the Green--Kubo relation
gives:
\begin{equation}
  \Dth = \frac{\kB T}{m\gamma_0} .
  \label{eq:GreenKubo}
\end{equation}

\subsection{S1.4. Nelson diffusion}

In Nelson's stochastic mechanics, the diffusion coefficient
associated with the forward and backward Wiener processes is:
\begin{equation}
  \Dq = \frac{\hbar}{2m} .
  \label{eq:Nelson}
\end{equation}

\subsection{S1.5. Closure to theorem}

Under Conjecture~1 ($\Dth = 2\Dq$):
\begin{align}
  \frac{\kB T}{m\gamma_0} &= 2\cdot\frac{\hbar}{2m}
  = \frac{\hbar}{m} \\
  \Rightarrow\quad \gamma_0 &= \frac{\kB T}{\hbar} .
  \label{eq:closure}
\end{align}
Then $R = \Dth/\Dq = 2$ and
$\Ebr = \kB T\ln 2 = \EL$.

% ============================================================================
\section{S2. Three Kernel Solutions and Crossing Times}
\label{sec:kernels}
% ============================================================================

We give the effective friction, TICE closure factor, and
crossing time $t^*$ for three physically motivated kernels.

\subsection{S2.1. Ohmic--Drude (exponential) kernel}

$K(\tau) = A\,e^{-\Omega|\tau|}$, with
$A = m\kB T\gamma_0\Omega/\hbar$.

Effective friction integrated against $|\psi|^2 = 1$
(uniform profile):
\begin{align}
  \geff(t) &= \gamma_0 + \frac{\beta A}{\Omega}
  \bigl(1 - e^{-\Omega t}\bigr) , \\
  \CTICE(t) &= \frac{\kB T/\hbar}
  {\gamma_0 + (\beta A/\Omega)(1 - e^{-\Omega t})} , \\
  t^*_{\mathrm{Drude}} &= \frac{1}{\Omega}\ln\!\Biggl(
  \frac{\beta A/\Omega}
  {\gamma_0 + \beta A/\Omega - \kB T/\hbar}\Biggr) .
  \label{eq:tstar_Drude}
\end{align}

\textit{Existence condition.}
$t^*$ is real and positive iff
$\gamma_0 < \kB T/\hbar < \gamma_0 + \beta A/\Omega$.

\subsection{S2.2. Algebraic (power-law) kernel}

$K(\tau) = A\,(1 + \Omega|\tau|)^{-n}$ with $n > 1$.

Integrated friction:
\begin{equation}
  \geff(t) = \gamma_0 + \frac{\beta A}{\Omega(n-1)}
  \Bigl[1 - (1+\Omega t)^{1-n}\Bigr] .
\end{equation}
Asymptote: $\geff(\infty) = \gamma_0 + \beta A/[\Omega(n-1)]$.

Crossing time (implicit):
\begin{equation}
  (1 + \Omega t^*)^{n-1}
  = \frac{\beta A/[\Omega(n-1)]}
  {\gamma_0 + \beta A/[\Omega(n-1)] - \kB T/\hbar} .
  \label{eq:tstar_alg}
\end{equation}

\subsection{S2.3. Gaussian kernel}

$K(\tau) = A\,e^{-\Omega^2\tau^2/2}$.

Integrated friction:
\begin{equation}
  \geff(t) = \gamma_0 + \beta A\sqrt{\frac{\pi}{2\Omega^2}}
  \,\mathrm{erf}\!\Bigl(\frac{\Omega t}{\sqrt{2}}\Bigr) .
\end{equation}
Asymptote:
$\geff(\infty) = \gamma_0 + \beta A\sqrt{\pi/(2\Omega^2)}$.

Crossing time (implicit):
\begin{equation}
  \mathrm{erf}\!\Bigl(\frac{\Omega t^*}{\sqrt{2}}\Bigr)
  = \frac{\kB T/\hbar - \gamma_0}
  {\beta A\sqrt{\pi/(2\Omega^2)}} .
  \label{eq:tstar_gauss}
\end{equation}

% ============================================================================
\section{S3. Perturbative Expansion of Deviation Law}
\label{sec:perturbation}
% ============================================================================

Let $\gamma_0 = \gamma^*(1 + \delta)$ with
$\gamma^* = \kB T/\hbar$.  The deviation from Landauer is:
\begin{equation}
  \Delta E = \Ebr - \EL = -\kB T\ln(1+\delta) .
  \label{eq:dev_exact}
\end{equation}

Taylor expansion for $|\delta| \ll 1$:
\begin{equation}
  \Delta E = -\kB T \Bigl[\delta
  - \frac{\delta^2}{2}
  + \frac{\delta^3}{3} - \cdots\Bigr] .
  \label{eq:dev_taylor}
\end{equation}

For the TICE time-dependent case with $\CTICE(t) = 1 + \epsilon(t)$:
\begin{equation}
  \Delta E(t) = \kB T\ln(1+\epsilon)
  \approx \kB T\Bigl[\epsilon
  - \frac{\epsilon^2}{2} + \cdots\Bigr] .
  \label{eq:dev_TICE}
\end{equation}

The sign convention: $\epsilon > 0 \Rightarrow \CTICE > 1 \Rightarrow$
super-Landauer (reversible work extraction possible).
$\epsilon < 0 \Rightarrow$ sub-Landauer (second-law violation for
irreversible erasure).

% ============================================================================
\section{S4. Nelson/TICE Internal Closure Ratio}
\label{sec:ratio}
% ============================================================================

Define the Nelson--TICE internal ratio:
\begin{equation}
  \mathcal{R}(t) \equiv
  \frac{\text{thermal diffusion at time }t}
  {\text{Nelson diffusion}}
  = \frac{\kB T}{m\,\geff(t)} \cdot \frac{2m}{\hbar}
  = 2\,\CTICE(t) .
  \label{eq:Rinternal}
\end{equation}

Timeline:
\begin{itemize}
  \item $t = 0$: $\mathcal{R}(0) = 2\kB T/(\hbar\gamma_0)$.
  If $\gamma_0 > \kB T/\hbar$ (overdamped initial condition)
  then $\mathcal{R}(0) < 2$.
  \item $t = t^*$ (crossing): $\mathcal{R}(t^*) = 2$,
  exactly Landauer.
  \item $t \to \infty$: $\mathcal{R}(\infty) =
  2\kB T/[\hbar\,\geff(\infty)]$.
\end{itemize}

The Landauer crossing is a \emph{transient event} in the TICE
framework, occurring when the kernel-accumulated friction
drives $\geff$ through $\kB T/\hbar$.

% ============================================================================
\section{S5. Bochner Positivity: What It Does and Does Not Prove}
\label{sec:bochner}
% ============================================================================

Bochner's theorem states that a function $K:\mathbb{R}\to\mathbb{C}$
is positive-definite (i.e., a valid bath correlation function) if
and only if it is the Fourier transform of a non-negative measure.

\textit{What Bochner guarantees:}
All three kernels in Sec.~\ref{sec:kernels} (Drude, algebraic,
Gaussian) are positive-definite for $A, \Omega > 0$ and hence are
admissible bath correlators.

\textit{What Bochner does not guarantee:}
Admissibility does not imply \emph{uniqueness}.
Infinitely many positive-definite kernels exist.
The closure $\Dth = 2\Dq$ selects a relationship between
$\gamma_0$ and $T$ but does not select a unique kernel.
The kernel shape determines how $\geff(t)$ approaches the
closure, but the asymptotic closure ratio is the same for
all kernels once the integral of $K$ is fixed.

% ============================================================================
\section{S6. Numerical Verification Summary}
\label{sec:numerical}
% ============================================================================

All results were verified numerically by the Q-Bond Research
Engine (15 audit scripts, 388 checks, 0 failures).
BT\#11 specifically was tested in scripts
\texttt{bt11\_verification.py} and \texttt{bt19\_verification.py}
with the following checks:

\begin{center}
\begin{tabular}{lll}
\hline
Test & Expected & Status \\
\hline
$R = 2$ at $\gamma_0 = \kB T/\hbar$ & $2.000$ & PASS \\
$\Ebr/(\kB T) = \ln 2$ & $0.6931\ldots$ & PASS \\
$\gamma_0\hbar/(\kB T) = 1$ & $1.000$ & PASS \\
$\tcor = \hbar/(\kB T)$ & exact match & PASS \\
$\partial R/\partial T = 0$ & $0.000$ & PASS \\
Deviation: $\delta = 0.1$ & $\Delta E/\kB T = -0.0953$ & PASS \\
Deviation: $\delta = -0.1$ & $\Delta E/\kB T = +0.1054$ & PASS \\
Sign rule: $\delta > 0 \Rightarrow \Delta E < 0$ & correct & PASS \\
Sign rule: $\delta < 0 \Rightarrow \Delta E > 0$ & correct & PASS \\
Mass independence of $R$ & verified 5 masses & PASS \\
Ohmic kernel crossing & $t^* > 0$, real & PASS \\
$\CTICE(t^*) = 1$ at crossing & $1.000$ & PASS \\
Gaussian kernel crossing & $t^* > 0$, real & PASS \\
Algebraic kernel crossing & $t^* > 0$, real & PASS \\
\hline
\end{tabular}
\end{center}

All numerical checks use double precision with tolerance
$10^{-12}$.

% ============================================================================
\section{S7. Consistency with Known Results}
\label{sec:consistency}
% ============================================================================

\subsection{S7.1. Caldeira--Leggett limit}

In the Markov limit ($\Omega \to \infty$), the memory kernel
$K(\tau) \to 2A\,\delta(\tau)/\Omega$ and $\geff \to \gamma_0$
instantaneously.  The TICE closure factor reduces to its Markov
value $\CTICE = \kB T/(\hbar\gamma_0)$, recovering the main
theorem.

\subsection{S7.2. Fluctuation--dissipation theorem}

At $\CTICE = 1$, the relation $\geff = \kB T/\hbar$ implies
$\Dth = \kB T/(m\geff) = \hbar/m = 2\Dq$, confirming
self-consistency with the FDT at the closure point.

\subsection{S7.3. Landauer bound}

The second law requires $\Ebr \geq \EL$ for irreversible
one-bit erasure.  The deviation formula gives
$\Ebr - \EL = -\kB T\ln(1+\delta) \geq 0$ iff
$\delta \leq 0$ (underdamped), consistent with the
physical requirement that reaching Landauer requires
approaching the closure from below.

\end{document}
